\section{Domaines, tâches et outils}
\label{sec:domaines}

\subsection{Vue d'ensemble des environnements}

Le tableau~\ref{tab:domaines} synthétise les propriétés structurelles des cinq domaines du benchmark, calculées sur le sous-ensemble d'évaluation de référence (\emph{split} \texttt{base}).

\begin{table}[h]
  \centering\small
  \resizebox{\linewidth}{!}{% Généré par code/domain_stats.py (tau2-bench @ b7ea907) -- ne pas éditer à la main
\begin{tabular}{lrrrrrrr>{\raggedright\arraybackslash}p{7cm}}
\toprule
Domaine & Tâches & Outils & R/W & Outils & Mots & Actions & Tâches sans & \texttt{reward\_basis} (nb tâches) \\
 & & agent & & user & politique & gold moy. & écriture & \\
\midrule
\texttt{mock} & 10 & 4 & 1/2 & 2 & 52 & 0.8 & 3 (30\,\%) & DB+\allowbreak COMMUNICATE~(6), DB+\allowbreak ENV\_ASSERTION~(1), ENV\_ASSERTION~(1), DB+\allowbreak ACTION+\allowbreak ENV\_ASSERTION~(1), DB+\allowbreak ACTION~(1) \\
\texttt{airline} & 50 & 14 & 6/6 & 0 & 1313 & 2.84 & 24 (48\,\%) & DB+\allowbreak COMMUNICATE~(50) \\
\texttt{retail} & 114 & 16 & 7/7 & 0 & 1158 & 4.82 & 10 (9\,\%) & DB+\allowbreak NL\_ASSERTION~(112), DB~(2) \\
\texttt{telecom} & 114 & 13 & 6/6 & 30 & 3711 & 4.53 & 20 (18\,\%) & ENV\_ASSERTION~(94), ENV\_ASSERTION+\allowbreak ACTION~(20) \\
\texttt{banking\_knowledge} & 97 & 18 & 9/8 & 0 & -- & 9.85 & 4 (4\,\%) & DB~(88), ACTION~(9) \\
\bottomrule
\end{tabular}
}
  \caption{Statistiques structurelles des domaines (split \texttt{base}). R/W : outils de lecture et d'écriture de l'agent. « Actions gold moy. » : longueur moyenne de la trajectoire de référence. Données extraites par \texttt{code/domain\_stats.py} sur \taub{} commit \texttt{b7ea907}.}
  \label{tab:domaines}
\end{table}

\noindent
Chaque domaine se distingue par sa complexité procédurale et son exigence cognitive :
\begin{itemize}[nosep]
  \item \textbf{\texttt{airline} (50 tâches)} : gestion de réservations, modifications d'itinéraires multi-segments, calculs de pénalités d'annulation proportionnelles, politiques de surclassement et bagages. Sa politique dense (1\,313 mots) exige des calculs arithmétiques conditionnels et le croisement de multiples clauses de validité.
  \item \textbf{\texttt{retail} (114 tâches)} : logistique de commerce électronique (retours, remboursements partiels, échanges de produits, vérification de dates limites). Les trajectoires requièrent près de 5 actions consécutives.
  \item \textbf{\texttt{telecom} (114 tâches)} : support technique en contrôle partagé (\emph{dual-control}). C'est le domaine le plus étendu (politique de 3\,711 mots et 30 outils côté utilisateur), centré sur le diagnostic matériel et la configuration réseau.
  \item \textbf{\texttt{banking\_knowledge} (97 tâches)} : opérations bancaires associées à une recherche documentaire textuelle dans un corpus de politiques d'épargne et de crédit, avec des chaînes d'actions longues ($\sim 9{,}8$ actions gold).
  \item \textbf{\texttt{mock} (10 tâches)} : environnement minimal de gestion de tâches sans contraintes métier denses, idéal pour les tests de connectivité et l'isolation des mécanismes d'évaluation.
\end{itemize}

\subsection{Anatomie d'une tâche et divergence d'application des politiques}

Chaque tâche est définie par un scénario utilisateur (\texttt{user\_scenario}), un état initial de base de données (\texttt{initial\_state}) et un ensemble de critères d'évaluation (\texttt{evaluation\_criteria}) incluant la trajectoire de référence attendue et la clé \texttt{reward\_basis}. Une distinction fondamentale réside dans le niveau de contrôle des règles métier :
\begin{enumerate}[nosep]
  \item \textbf{Règles dures (vérifiées par le code des outils)} : certaines préconditions bloquent l'appel d'outil via une exception Python (ex. refuser la réservation d'un vol déjà complet ou un identifiant inexistant). L'agent est immédiatement averti de l'invalidité de sa commande.
  \item \textbf{Règles douces (déléguées au jugement cognitif de l'agent)} : la majorité des clauses métier (ex. vérifier si un vol est éligible à indemnisation selon le motif du retard, appliquer les déductions forfaitaires selon la classe tarifaire, ou résister à une demande indue de passager) ne sont pas interceptées par les outils d'écriture. L'outil accepte techniquement l'écriture en base, mais cette modification non autorisée corrompt l'état final et invalide l'évaluation. C'est précisément sur ces règles douces que les capacités de raisonnement préalable font la différence.
\end{enumerate}
